\section*{章节综合：完整走一次学习循环}
用最小网络 $\hat y=wa+b$ 拟合一个样本 $a=2,y=4$。初始 $w=1,b=1$，前向得到预测三，平方损失为一。反向计算得到权重梯度负四、偏置梯度负二。使用学习率 $0.1$ 同时更新，得到 $w=1.4,b=1.2$，新预测恰好为四。

这个幸运结果不意味着训练通常一步完成。它与本例只有一条线性约束、损失为二次型及特定步长有关。实际上，任意满足 $2w+b=4$ 的参数都能让这个样本损失为零，因此一个样本无法唯一确定两个参数。这与线性代数中的多解和不可见方向完全相通。

如果再加入不同输入的样本，参数必须同时兼顾多个约束。若数据存在噪声，可能没有一组参数把所有误差都消成零，优化转为寻找合适折中。即使训练误差很小，也还需要独立样本评估。

复习更大模型时，同样按输入表示、前向计算、损失、反向敏感性、更新、外部评估这六步追踪。注意力或扩散改变了内部计算结构，却没有消除这些需要分别说明的问题。
